\documentclass[a4paper, 12pt]{article}

\usepackage[utf8]{inputenc}
\usepackage[T1]{fontenc}
\usepackage[spanish, es-tabla]{babel}
\usepackage{lmodern}
\usepackage[top=2.5cm, bottom=2.5cm, left=2.5cm, right=2.5cm]{geometry}
\usepackage{graphicx}
\graphicspath{{Figures/}}
\usepackage[dvipsnames]{xcolor}
\usepackage{microtype}

\definecolor{uni-blue}{RGB}{0, 56, 101}
\definecolor{uni-red}{RGB}{166, 25, 46}

\pagestyle{empty}

\begin{document}

\begin{titlepage}
    \centering
    \vspace*{0.4cm}
    
    %% ENCABEZADO INSTITUCIONAL
    {\Large\bfseries UNIVERSIDAD NACIONAL DE INGENIER\'IA}\\[0.35cm]
    {\large\bfseries FACULTAD DE CIENCIAS}\\[0.2cm]
    {\large Escuela de F\'isica}\\[1.0cm]
    
    %% ESCUDO OFICIAL UNI
    \includegraphics[width=4.3cm, keepaspectratio=true]{uni_logo.jpg}\\[1.0cm]
    
    %% CURSO Y MESA
    {\large\bfseries Física General / Física III (CF1B2) --- Sección A}\\[0.25cm]
    {\normalsize\textbf{Semestre Académico 2026-II \quad $\vert$ \quad Mesa de Trabajo N.° 3}}\\[0.8cm]
    
    %% TÍTULO DEL LABORATORIO
    {\large\textbf{LABORATORIO N.° 1:}}\\[0.35cm]
    {\LARGE\bfseries El osciloscopio como instrumento de medida\par}\vspace{0.9cm}
    
    %% INTEGRANTES
    \begin{minipage}{0.85\linewidth}
        \textbf{\large INTEGRANTES:}\\[0.35cm]
        \begin{tabular}{@{\hspace{0.5cm}}ll}
            \textbf{1. Vera Vivanco, Jesús Rodolfo} & (Código: \textbf{20252703K}) \\[0.25cm]
            \textbf{2. Casañe Ramírez, Gerald Giordano} & (Código: \textbf{20251368C}) \\[0.25cm]
            \textbf{3. Sánchez Feril, Piero Alejandro} & (Código: \textbf{20251629A})
        \end{tabular}
    \end{minipage}
    
    \vfill
    
    %% DOCENTE Y FECHA
    \begin{minipage}{0.85\linewidth}
        \centering
        \textbf{\normalsize Docente:} Prof. Cortez R., Gregorio\\[0.25cm]
        \textbf{\normalsize Fecha de Entrega:} Lima, 11 de septiembre de 2026
    \end{minipage}
    \vspace*{0.4cm}

\end{titlepage}

\end{document}
